% =====================================================================
% chapters/01_Introduction.tex — rendu LaTeX de ../traductions/01_introduction.en_brut.txt (23 septembre 2026)
% Section 1 de la VERSION FRANÇAISE de l'article court. \input depuis main.tex, première du corps.
% Ce fichier DÉFINIT \label{sec:intro}. La section 1.4 renvoie à toutes les autres sections :
%   elle porte six \ref.
% Citations : lammer2026gta (GTA), vu2025modeling (architecture), tisseo2023emc2 (enquête).
% Rendu mot pour mot de la traduction fournie par l'auteur.
% =====================================================================

\section{Introduction}
\label{sec:intro}

\subsection{Les promesses des agents génératifs}
\label{sec:promise}

Les simulations multi-agents existantes de mobilité urbaine multimodale reposent sur des modèles tabulaires estimés à partir d'enquêtes de mobilité des ménages, ou sur des règles explicites définies par des experts du domaine. Une enquête de mobilité des ménages demande aux habitants d'un territoire de décrire chaque déplacement effectué au cours d'une journée. Dans les deux cas, l'espace des comportements représentables est fixé \emph{a priori} ; un facteur n'étant ni une variable ni une règle ne peut influencer aucune décision.

Cet espace comportemental fixe limite précisément le domaine où la simulation est censée intervenir. Une ville teste une politique de transport sur une population simulée avant d'entreprendre toute construction. Le résultat obtenu est limité par le comportement que son modèle peut représenter. Le choix du mode de transport étant multidimensionnel, intégrer cette complexité à l'échelle de la ville est difficile.

Les agents génératifs construits sur des modèles de langage promettent de surmonter cette difficulté. Alimentés par des corpus massifs, ils intègrent des heuristiques de décision que les enquêtes n'enregistrent pas et s'adaptent sans règles explicites. Une grève, une vague de chaleur ou un article de presse sur la sécurité dans le métro pourraient alors aboutir à une décision qu'aucune enquête ne permet de prendre. Ces systèmes promettent également de fonctionner sans les données d'étalonnage locales requises par les modèles à base d'agents. Des ensembles de données détaillés sur les trajectoires multimodales existent uniquement pour quelques grandes métropoles.

\subsection{Évaluation des promesses face à une population réelle}
\label{sec:gap}

Les évaluations publiées de ces agents font rarement l'objet d'une comparaison avec une population humaine. La plupart évaluent un agent génératif par rapport à un agent basé sur des règles, ou par rapport à des modèles agrégés produits par les auteurs eux-mêmes. Un système se distingue.

GTA \citep{lammer2026gta} confronte une répartition modale simulée à une enquête nationale sur les déplacements. La répartition modale correspond à la part des déplacements effectués par chaque mode. GTA crée ses personas à partir de cette enquête. Un modèle de langage rédige le plan de journée de chacun. Sa seule référence modélisée est la répartition observée dans d'autres régions. Copier les parts de Hambourg génère une erreur quadratique moyenne de 1,99 sur Berlin, contre 4,07 pour sa propre simulation.

Cet article apporte quatre éléments manquants à ces comparaisons. Une ligne de base sous les agents indique la décision qu'un décideur prendrait sans connaissance comportementale du territoire. Les modèles ajustés sur l'enquête propre à ce territoire indiquent ce que confirment ses données. Les mêmes entrées des deux côtés créent un écart attribuable au décideur plutôt qu'aux informations qui lui ont été communiquées. Une lecture sous l'agrégat indique si deux décideurs produisant les mêmes parts prennent des décisions similaires.

\subsection{Ce que fait cet article}
\label{sec:contributions}

Cet article mesure la pertinence de la délibération verbalisée dans un agent de mobilité. La délibération verbalisée consiste pour le modèle à formaliser son raisonnement par écrit avant de répondre. Dans une population d'agents de mobilité, elle n'améliore pas la fidélité au régime ordinaire décrit par l'enquête. Elle introduit des événements qu'aucune variable n'encode. Nous retraçons le parcours d'un tel événement jusqu'à la décision.

Nous effectuons des mesures sur un territoire, la région toulousaine, en nous basant sur son enquête certifiée sur les déplacements des ménages de 2023 \citep{tisseo2023emc2}. Un groupe contrôlé de 1\,000 profils synthétiques effectue les 3\,299 déplacements du jour que nous analysons. Quinze décideurs interviennent sur ces déplacements. Nous appelons «~décideur~» tout agent qui transforme la description d'un déplacement en une probabilité parmi les options proposées. Les agents évoluent dans une simulation multi-agents GAMA de la région toulousaine, avec ses réseaux et horaires réels. Cette simulation suit l'architecture GAMA--OpenTripPlanner--LLM de \citet{vu2025modeling}.

Cet article repose sur trois contributions. La première est un banc d'essai comparatif avec des entrées égales entre agents génératifs et modèles tabulaires, sur une cohorte synthétique contrôlée par rapport à cette enquête. La seconde concerne la position de quinze décideurs sur ce banc d'essai, et deux dissociations que les parts agrégées ne révèlent pas. Le troisième point concerne le parcours d'un événement non tabulé jusqu'à la décision, au sein d'un agent génératif doté de mémoire.

\subsection{Organisation}
\label{sec:organisation}

Cet article suit l'ordre de la mesure. La section~\ref{sec:related} le situe dans trois contextes bibliographiques. La section~\ref{sec:agent} décrit ce que l'agent évalué reçoit et renvoie. La section~\ref{sec:bench} présente le banc de comparaison et son protocole. La section~\ref{sec:results} décrit une journée type. La section~\ref{sec:nontabulated} analyse l'événement non tabulé. La section~\ref{sec:implications} en tire les implications, les limites et la conclusion.
